% Chapter 8 -- The command-line compiler (model: TP 6.0 UG Ch 9).  The
% reference for TPS.COM: its files, every line it prints (KIND protocol and
% dumprow, registered against the generated inventory both ways), the exit
% code, and every capacity with its value as shipped (KIND limit).
% Sources (Law 1): os/tps/tps.pas (the driver, RunLexer/RunParser/
% RunTypeChecker/RunCodeGenerator, DumpToken, TraceEnter/TraceExit, the
% symbol dump writers, the constants SourceMax ... PendingNameMax and the
% checks that use them); os/tps/LEXER-DUMP.md, PARSER-TRACE.md,
% SYMBOLS-DUMP.md, CODEGEN.md s1; seed/rt/fileio.asm (missing TPSIN.PAS).
% The listings are replayed by test-pubs-turbo; test-tps-*-os prove TPS.COM
% prints the same on InitechDOS.
\Gates{test-tps-lex-fpc, test-tps-lex-os, test-tps-parse-fpc, test-tps-parse-os, test-tps-type-fpc, test-tps-type-os, test-tps-gen-fpc, test-tps-gen-os, test-pubs-turbo}
\chapter{The command-line compiler}
\ix{command-line compiler}
\ix{TPS command}

This chapter is the reference for the compiler command, \file{TPS.COM}: the
files it uses, everything it prints, and its capacities.

\section{The TPS command}
\SyntaxLine{TPS}

\file{TPS.COM} takes no parameters and has no options.  It reads the
program from the file \file{TPSIN.PAS} in the current directory and, when
the program compiles, writes the assembly language to \file{TPSOUT.S} in
the same directory.  Everything else it reports goes to COM1.

If \file{TPSIN.PAS} does not exist, \file{TPS.COM} stops with the message
\code{FILEIO-ERROR AX=2} on COM1, as every Turbo Initech program does when a
file it opens is missing (the \doctitle{Programmer's Guide}, Chapter 12).
\file{TPS.COM} ends with exit code 0 whenever it has read the whole program,
whether or not the program compiled; the lines it prints say whether it
did.  \file{TPSOUT.S} is written only by a compilation in which all four
passes succeed.

\section{The output}
\ix{compiler output}
The compiler reads the program four times.  The report of each pass
begins and ends with a line of its own, and either reports success or one
error:

\begin{BorTable}{Lines of the compiler's output}{p{36mm}p{62mm}}
\Hd{Line} & \Hd{Meaning}\EndHd
\ProtocolItem{TPS-LEX-BEGIN} & start of the token listing\\
\ProtocolItem{TPS-LEX-ERROR} & a lexical error, with its position (ends the
  listing)\\
\ProtocolItem{TPS-LEX-END} & end of the token listing\\
\ProtocolItem{TPS-PARSE-BEGIN} & start of the syntax trace\\
\ProtocolItem{TPS-PARSE-OK} & the program follows the grammar\\
\ProtocolItem{TPS-PARSE-ERROR} & a syntax error\\
\ProtocolItem{TPS-PARSE-SKIP} & the pass was skipped: \code{lexer-error}\\
\ProtocolItem{TPS-PARSE-END} & end of the syntax trace\\
\ProtocolItem{TPS-TYPE-BEGIN} & start of the symbol listing\\
\ProtocolItem{TPS-TYPE-OK} & the declarations and types are correct\\
\ProtocolItem{TPS-TYPE-ERROR} & a semantic error\\
\ProtocolItem{TPS-TYPE-SKIP} & skipped: \code{lexer-error} or
  \code{parser-error}\\
\ProtocolItem{TPS-TYPE-END} & end of the symbol listing\\
\ProtocolItem{TPS-GEN-BEGIN} & start of code generation\\
\ProtocolItem{TPS-GEN-OK} & \file{TPSOUT.S} written: \code{bytes=} its
  size, \code{labels=} the number of labels\\
\ProtocolItem{TPS-GEN-ERROR} & a code-generation error\\
\ProtocolItem{TPS-GEN-SKIP} & skipped: \code{lexer-error},
  \code{parser-error} or \code{type-error}\\
\ProtocolItem{TPS-GEN-END} & end of code generation\\
\end{BorTable}

The error lines and their messages are described in the
\doctitle{Programmer's Guide}, Appendix A.

\section{The token listing}
\ix{token listing}
Between \code{TPS-LEX-BEGIN} and \code{TPS-LEX-END} the compiler lists one
token per line:

\begin{BorTable}{Rows of the token listing}{p{36mm}p{62mm}}
\Hd{Row} & \Hd{Token}\EndHd
\DumpRowItem{KW} \textit{word} & a reserved word, in small letters\\
\DumpRowItem{IDENT} \textit{name} & an identifier, in small letters\\
\DumpRowItem{INTEGER} \textit{value} & a number\\
\DumpRowItem{STRING} \textit{n} \textit{c}\dots & a string literal: its
  length, then the code of each character\\
\code{SEMI}, \code{ASSIGN}, \dots & a symbol, by its kind name
  (\doctitle{Programmer's Guide}, Appendix B)\\
\code{EOF} & the end of the file\\
\end{BorTable}

Comments and blanks produce no rows.  Writing each character of a string by
its code makes apostrophes and blanks unambiguous:

\Program{strs}

\CompilerOutput{strs}{LEX}

\section{The syntax trace}
\ix{syntax trace}
Between \code{TPS-PARSE-BEGIN} and \code{TPS-PARSE-END} the compiler writes
a row each time it begins and finishes recognizing a rule of the grammar:

\begin{BorTable}{Rows of the syntax trace}{p{36mm}p{62mm}}
\Hd{Row} & \Hd{Meaning}\EndHd
\DumpRowItem{ENTER} \textit{rule} \code{L}\textit{n} & the rule begins; the
  next token is on line \textit{n}\\
\DumpRowItem{EXIT} \textit{rule} \code{L}\textit{n} & the rule has been
  recognized; the token after it is on line \textit{n}\\
\end{BorTable}

The rules are named as the syntax diagrams of the \doctitle{Programmer's
Guide} are.  A rule in which an error is found has an \code{ENTER} row and no
\code{EXIT} row.  Chapter 4 shows how to use the trace.

\section{The symbol listing}
\ix{symbol listing}
Between \code{TPS-TYPE-BEGIN} and \code{TPS-TYPE-END} the compiler lists
every name the program declares, global names first in the order of their
declarations, then the names of each routine:

\begin{BorTable}{Rows of the symbol listing}{p{36mm}p{62mm}}
\Hd{Row} & \Hd{Meaning}\EndHd
\DumpRowItem{SCOPE} \code{global bytes=}\textit{n} & the global names;
  \textit{n} bytes of global variables\\
\code{SCOPE routine name=}\textit{r} \code{frame=}\textit{n} & the names of
  routine \textit{r}; a frame of \textit{n} bytes\\
\DumpRowItem{SYM} \code{kind=}\textit{k} \code{name=}\textit{x}
  \code{type=}\textit{t} \dots & one name: its kind, name and type, then the
  fields for the kind\\
\DumpRowItem{FIELD} \code{owner=}\textit{r} \code{index=}\textit{i} \dots &
  a field of record type \textit{r}, with its type, offset and size\\
\end{BorTable}

The kinds are \code{const}, \code{type}, \code{var}, \code{procedure},
\code{function}, \code{value-param}, \code{var-param}, \code{result} and
\code{local}.  Constants show their \code{value=}, or \code{value-length=}
for a string; variables, parameters, results and locals their
\code{offset=} and \code{size=} in bytes, and \code{array=}\textit{lo}..\textit{hi}
for an array; routines the number of \code{params=} and their \code{frame=};
record types the number of \code{fields=} and their \code{size=}.  The
\doctitle{Programmer's Guide}, Chapter 10, explains the offsets.

\Program{kinds}

\CompilerOutput{kinds}{TYPE}

\ProgramOutput{kinds}

\section{Capacities}
\ix{capacities}
\ix{limits, compiler}
The compiler keeps every table in a fixed array.  These are its
capacities:

\begin{BorTable}{Compiler capacities}{p{30mm}p{14mm}p{52mm}}
\Hd{Capacity} & \Hd{Value} & \Hd{Limits}\EndHd
\LimitItem{SourceMax}{262,144} & 262,144 & characters in the source file\\
\LimitItem{PoolMax}{6,144} & 6,144 & characters in the spellings of all
  the different identifiers together\\
\LimitItem{NameMax}{704} & 704 & different identifiers\\
\LimitItem{SymbolMax}{1,024} & 1,024 & rows of the symbol table: every
  constant, type, field, variable, routine, parameter, result and local\\
\LimitItem{HeaderParamMax}{32} & 32 & parameters of a routine; arguments of
  a call\\
\LimitItem{PendingNameMax}{32} & 32 & names in one declaration list, such
  as \code{var A, B, C: integer}\\
\LimitItem{ChunkMax}{255} & 255 & characters the compiler reads from the
  source at one time; it is not a limit on lines\\
\end{BorTable}

A program beyond the first three capacities is reported as a lexical error
(\code{input exceeds SOURCEMAX=}, \code{char pool exceeds POOLMAX=},
\code{name table exceeds NAMEMAX=}); beyond the next three, as
\code{semantic rule violation}.  In addition, an identifier or a string
literal may have at most 255 characters.
